\section{平稳分布、可逆性与遍历收敛}
\label{sec:06-Random-Processes/03-Markov-Chains/04}

你手上有一条 Markov 链，\(P\) 已经写好了。盯着转移矩阵看了一会，一个念头冒出来：这链跑久了，状态分布会不会停下来？

把你猜测的那个``停下来''的分布叫做 \(\pi\)。如果它一步之后纹丝不动：

\[
\pi P=\pi,
\qquad
\sum_i\pi_i=1,\quad\pi_i\ge0,
\]

就叫\textbf{平稳分布}。只要 \(X_0\sim\pi\)，每个 \(X_n\) 的边缘分布都还是 \(\pi\)，因为每步转移就是把分布右乘一次 \(P\)。

注意``每个时刻的边缘分布相同''和``路径不动''是两码事。链仍然在状态之间跳来跳去，只不过从任何一个状态流出的概率总量，恰好被从其他状态流入的概率总量补上了。好比一个湖泊，蒸发和降水各自稳定，但每一滴水并不停在原处。

\subsection{同一个平稳分布，三种不同的收敛}

你现在觉得事情很简单：只要解出 \(\pi P=\pi\)，链就一定收敛过去？试试看。

\hyperref[sec:06-Random-Processes/03-Markov-Chains/02]{``可达、通信类与周期性''}一节里那个两状态确定交替链：\(P_{12}=P_{21}=1\)，每个状态隔一步必须换一次。它的平稳分布是 \((1/2,1/2)\)。但如果你从状态一出发，时刻分布永远是

\[
(1,0),(0,1),(1,0),(0,1),\ldots
\]

在 \((1,0)\) 和 \((0,1)\) 之间来回跳，根本不收敛。

这个例子把三个容易混在一起的概念拆开了：

\begin{enumerate}
\tightlist
\item
  \textbf{平稳分布}：固定点方程 \(\pi P=\pi\) 的解。它存在，但链不一定往那里收敛；
\item
  \textbf{极限分布}：对任意初始分布 \(\mu_0\)，\(\mu_0P^n\) 是否当 \(n\to\infty\) 时趋于某个确定的分布？上面那个交替链的答案是否；
\item
  \textbf{时间占比}：在一条具体的长路径里，访问每个状态的频率趋于多少？交替链里，每个状态恰好占一半时间，和稳态一致。
\end{enumerate}

三个概念在适当的遍历条件下会重合，但定义上互不隶属。你需要的答案类型决定了你应该关心哪一个——MCMC 采样关心的是时间平均，做短期预测可能只关心极限分布，而分析平衡性质只需要平稳方程本身。

\subsection{有限不可约非周期链：三件事同时成立}

把条件收紧。对状态数有限的链：

\begin{itemize}
\item
  不可约 \(\Rightarrow\) 有唯一平稳分布 \(\pi\)，而且每个分量 \(\pi_i>0\)。唯一性来自 Perron--Frobenius：不可约非负矩阵的谱半径是单重正特征值，对应的正特征向量就是平稳分布。
\item
  再加上非周期 \(\Rightarrow\) 转移概率逐点收敛：
  \[
  P^n_{ij}\to\pi_j,
  \]
  对所有 \((i,j)\) 成立，不再依赖初始状态。
\item
  因此任意初始分布 \(\mu_0\) 都满足
  \[
  \mu_0P^n\to\pi.
  \]
\end{itemize}

每个条件负责一件事：不可约保证不存在多个互不通信的区域各自走向不同的平衡；正常返（有限链自动满足）保证平稳概率可以归一化；非周期消除节拍振荡——有了非周期，你不再需要跟踪``现在是奇数步还是偶数步''。

对无限状态链，正常返不是自动的，需要单独检验。有限链的便利在于它替你省掉了这些额外条件，但直觉不能照搬到无限情形。

\subsection{遍历定理：时间平均不需要非周期}

考虑一条不可约正常返链的一条长路径。对任意可积函数 \(f\)，在适当条件下：

\[
\frac1n\sum_{k=0}^{n-1}f(X_k)
\xrightarrow{\text{a.s.}}
\sum_j\pi_jf(j).
\]

取 \(f(j)=\mathbf1_{\{j=r\}}\)，得到状态 \(r\) 的长期访问比例：

\[
\frac1n\sum_{k=0}^{n-1}\mathbf1_{\{X_k=r\}}
\to\pi_r.
\]

这里没有要求非周期。周期链的时刻分布在振荡，但长期访问频率照样稳定。非周期性只负责逐时刻分布的收敛，即 \((P^n)_{ij}\to\pi_j\)；时间平均的收敛只要不可约加正常返就够。

这个区分是 MCMC 的命脉。你几乎不会靠等待某一时刻的分布刚好是 \(\pi\) 来做推断——你用一条长链的沿途平均值去估计 \(\pi\) 下的期望，需要的恰恰是时间平均收敛，而非逐时刻收敛。时间平均也不需要等到链``忘记''初始状态之后才开始记录，burn-in 只是减少初始偏差，不是理论前提。

\subsection{细致平衡：全局平衡的逐对版本}

平稳条件 \(\pi P=\pi\) 的每一行写开，是

\[
\sum_j\pi_jP_{ji}=\pi_i,
\]

流入 \(i\) 的总概率流等于流出 \(i\) 的总概率流。这是全局平衡。

现在加强一步：要求每一对 \((i,j)\) 之间概率流单独抵消：

\[
\pi_iP_{ij}
=\pi_jP_{ji}.
\]

这叫\textbf{细致平衡}。如果一个分布满足这个对所有 \((i,j)\) 成立，链相对于 \(\pi\) 称为\textbf{可逆}——因为时间反演后的转移概率和正向一样。

对 \(i\) 求和立刻得到平稳：

\[
\sum_i\pi_iP_{ij}
=\sum_i\pi_jP_{ji}
=\pi_j.
\]

所以细致平衡推出平稳，反之不必然。一个平稳链完全可以有环流：状态 A 流入 B，B 流入 C，C 流回 A，每一对之间的净流不为零，但每个状态的总流入仍然等于总流出。不可逆链在自然界普遍存在，比如带循环的化学反应网络和不可逆排队系统。

那为什么要关心细致平衡？因为它让构造、模拟和验证变得简单。Metropolis--Hastings 和 Gibbs 采样的核心策略正是：不是直接解 \(\pi P=\pi\)，而是逐对设计转移使得

\[
\pi_i\cdot P_{ij}
=\pi_j\cdot P_{ji}
\]

自动成立，从而保证目标分布是平稳分布。逐对验证比求解全局流方程组更机械，也更不容易遗漏。

\subsection{天气链：一个可以手指算出来的例子}

回到\hyperref[sec:06-Random-Processes/03-Markov-Chains/01]{``Markov 性质、转移矩阵与多步演化''}一节里的天气模型：

\[
P=
\begin{bmatrix}
0.8&0.2\\
0.4&0.6
\end{bmatrix}.
\]

令 \(\pi=(s,r)\)，晴天的平稳概率是 \(s\)，雨天是 \(r\)。平稳方程：

\[
0.8s+0.4r=s,
\]
等价于
\[
0.2s=0.4r,
\]

加上 \(s+r=1\)，得

\[
\pi=\left(\frac23,\frac13\right).
\]

这个等式 \(0.2s=0.4r\) 恰巧就是细致平衡：长期从晴转雨的流量

\[
\frac23\cdot0.2=\frac{2}{15},
\]

恰好等于从雨转晴的流量

\[
\frac13\cdot0.4=\frac{2}{15}.
\]

两状态链自动可逆，因为只有一对状态，环流无处可藏。链不可约（晴雨互达），且矩阵对角线 \(0.8\) 和 \(0.6\) 都大于零，所以非周期。结论：不论今天什么天气，日子足够长后，晴天的概率趋于 \(2/3\)，雨天趋于 \(1/3\)；一条很长的路径中，约三分之一的日子是雨天。

\subsection{收敛有多快：谱在背后工作}

平稳分布只回答``最终是什么''，不回答``多快到达''。对有限、不可约、可逆链，转移矩阵在以 \(\pi\) 加权的内积下是自伴的，因此特征值全是实数，可以排序：

\[
1=\lambda_1>\lambda_2\ge\cdots\ge\lambda_m\ge-1.
\]

\(\lambda_1=1\) 对应平稳分布本身，它不变。其余特征值控制偏离的衰减速度。令

\[
\lambda_*=\max_{i\ge2}|\lambda_i|.
\]

在合适的范数下，偏离平稳分布的主要成分以 \(\lambda_*^n\) 的速率衰减；相应的\textbf{绝对谱隙} \(1-\lambda_*\) 越大，混合通常越快。

只盯着 \(1-\lambda_2\) 会漏掉接近 \(-1\) 的特征值带来的麻烦——它们导致符号交替的振荡模态。不可约可逆链如果是周期的，则必为二周期，特征值 \(-1\) 出现；这时链不会收敛到平稳分布，而是在两个分布之间来回跳。懒化操作把转移矩阵换成 \(\frac{I+P}{2}\)，特征值映射为 \((1+\lambda_i)/2\)，把 \(-1\) 变成 \(0\)，消除振荡，同时保持平稳分布不变。

对不可逆链，转移矩阵在加权内积下不自伴，特征值可以为复数，矩阵的非正规性也会影响收敛行为，不能直接套用这组实特征值的说法。精确的混合时间还依赖于起始状态、状态空间的大小，以及你用什么距离度量来衡量``接近''——单凭``有唯一平稳分布''不能判断模拟是否已经跑够了。
